\section{Einzigartigkeit der Bibel}
Bevor wir anfangen, müssen wir uns zuerst gemeinsam ein paar Begriffe anschauen.
\begin{itemize}
    \item[Bibel:] Bibel kommt vom griechischen Wort \verbN[b, p]{biblios}. Was übersetzt einfach Buch heisst.
    \item[Altes Testament:] Das sind die Schrift die vor der Geburt Jesus geschrieben wurden. Es besteht aus XXX Büchern.
    \item[Neues Testament:] Geschrieben nach der Geburt und dem Tod von Jesus mit XXX Bücher. 
\end{itemize}
Das sind jetzt schon mal die Eckdaten. Ich habe hier vorne diverse Übersetzungen auf Deutsch. Schlachter, Elberfelder, Hoffnung für alle, Neue evangelistische Übesetzung, katholische Übersetzungen. Auch wenn es jeweils andere Übersetzungen sind, so haben doch alle den gleichen Inhalt. Die katholische Übesetzung hat noch die Apokryphen drin, was das ist werden wir uns noch anschauen.

\subsection{Details zur Bibel}
So hier ein paar Zahlen zur Bibel:
\begin{enumerate}
    \item Geschrieben über eine Zeitspanne von 1400 Jahren
    \item Geschrieben von mehr als 40 Authoren aus verschiedenen Gesellschaft schichten:
    \begin{itemize}
        \item Mose, in den Universitäten Ägyptens geschult;
        \item Petrus, ein Fischer;
        \item Amos, ein Hirte;
        \item Josua, ein Offizier;
        \item Nehemia, ein Mundschenk;
        \item Daniel, ein Ministerpräsident;
        \item Lukas, ein Arzt;
        \item Matthäus, ein Zollbeamter;
        \item Paulus, ein Rabbiner
    \end{itemize}
    \item Geschrieben an verschiedenen Orten:
    \begin{itemize}
        \item Mose in der Wüste;
        \item Jeremia im Kerker;
        \item Lukas, auf Reisen;
        \item Paulus, im Gefängnis;
        \item \dots
    \end{itemize}
    \item Geschrieben zu verschiedenen Zeitspanne
    \item Geschrieben in verschieden Gemütsverfassungen
    \item Geschrieben in drei Sprachen
    \begin{itemize}
        \item Hebräisch das Alte Testament
        \item Aramäisch Teile des Alten Testaments
        \item Griechisch das neue Testament
    \end{itemize}    
\end{enumerate}
Trotz dieser langer Zeit, den unterschiedlichen Autoren, Orte und Gegebenheiten, die die Bibel eine Einheit. Das Paradies in 1. Mose 1 taucht in Offenbarung dem letzten Buch der Bibel wieder auf. Dazwischen zieht der Heilsplan Gottes wie ein roter Faden durch die verschiedenen Bücher, mit dem ersten Höhepunkt vom ersten kommen Jesus, bis zum Ende beim zweiten kommen Jesus und Herstellung einer neuen Erde.

Es ist das einzige Buch auf der Welt, das konkrete Angaben für die Zunkunft gibt, sogenannte Prophetien. Und das nicht nur vage Angaben wie die Vorhersagen von Nostradamus, sondern konkrete nachweisbare Vorhersagen. So wird zum Beispiel im Buch Jeremia, der Namen vom persischen König Kyrus schon mehr als 100 Jahre vor seiner Geburt genannt. Oder die Phrophezeihung über das erste Kommen Jesu und über seinen Tod. 

\subsection{Sie hat die Zeit überlebt}
Wie schon gesagt, ist die Bibel über 1400 Jahre entstanden und ist inzwischen 3300 Jahre alt. Und das, obwohl sie auf vergänglichem Material geschrieben wurde und bis die Druckerpresse erfunden wurde immer abgeschrieben werden musste.

Es gibt 8000 Manuskripte der lateinischen Vulgata und wenigsten 1000 einer frühren Version. Dazu kommen 4000 griechische Manuskripte vom neuen Testament. Daneben kann ein grosser Teil des Neuen Testaments aus den Zitaten der frühchristlichen Verfasser neu erstellt werden.

So als Gegenüberstellung: Homer, Ilias 

\subsection{Sie hat Verfolgung überlebt}

Voltair, der 1778 gestorbene berühmte französiche Ungläugibe, behauptete, das Christentum werde einhundert Jahre nach seiner Zeit nur noch im Museum existieren. Tatsächlich nimmt jedoch die Verbreitung der Bibel zu.\\
Die Ironie der Geschichte, 60 Jahr nach seinem Tod übernahme die Bibelgeschellschaft Genf, sein Haus, seine Druckerpresse und druckte dort die Bibel.

Im Jahr 303 n. Chr. befahl der römische Kaiser Diokletian die bedingungslose Verfolgung der Christen. Eusebius schreibt in seiner Kirchengeschichte, dass ein kaiserlicher Erlass veröffentlicht wurde, welcher befahl, die Kirchen bis auf den Grund nieder zureißen und die Schriften zu verbrennen, und verfügte, dass Inhaber von
Ehrenstellen ihre Würden, und Bedienstete, sofern sie im Bekenntnis des Christentums verharrten, die Freiheit verlieren sollten.\\
25 Jahre später befahl Konstantin, der Nachfolger Diokletians, auf kosten der Regierung 50 Exemplare der Heiligen Schrift herszustellen.

\subsection{Sie hat Kritik überlebt}

Seit fast 2000 Jahren versuchen jetzt Ungläugibe dieses Buch die Bibel zu widerlegen. Aber die Bibel steht da wie ein Fels. Oft wurde schon der Untergang der Bibel prophezeit, aber wie wir hier sehen lebt sie immer noch in verschiedenen Farben und Formen.

Es passierte aber das Gegenteil. Viele Kritiker die Grossspurig verkündet haben, sie hätten die Bibel wiederlegt, wurden von der Bibel selber, aber auch von archäologischen Funden einen besseren belehrt und mussten sich beschämt verkriechen.

Als Beispiel die \betonung[b, p]{Quellenhypothese}. Sie besagt, weil es zur Zeit noch keine Schrift gab, wurden die 5 Bücher Moses später von verschieden Authoren geschrieben. Aber dann fand man die schwarze Stele. Sie war mit keilförmigen Schriftzeichen bedeckt und sie war um mindestens 3 Jahrhunderte älter als die Bücher Moses. Sie stammte also aus einer Zeit wo man glaubte, dass dort nur primitive Menschen ohne Alphabet gelebt hätten.

Apropo primitive Menschen, die Sprache des Menschen wurde im laufe der Zeit nicht komplexer sondern immer einfacher. Die alten Sprachen waren hochkomplexe Sprachen.

\subsection{Einzigartig in ihrer Geschichte}
Sie ist das einzige menschliche Werk, in dem sich eine Fülle von Prophetien in Bezug auf einzelne Nationen, auf Israel, auf alle Völker der Erde, auf bestimmte Stödte und auf das Kommen des Messias in erfüllung gingen. Keine anderes Buch, auch religöses Buch kann erfüllten Prophetien aufweisen. Auch der Islam hat keine Prophetie über das kommen von Mohamed.

Die Geschichte in diesem Buch ist historisch belegbar. Sie entspricht der Wahrheit. Oft wurde versucht dies zu wiederlegen, aber immer neue archäologische Funden haben diese Kritiker verstummen lassen.


